\section{Conclusiones principales}

El proyecto implement\'o y calibr\'o un sistema recomendador de inversiones para
FinPUC en tres capas metodol\'ogicas (Markowitz, Black-Litterman, Monte Carlo
P4). En esta fase, el objetivo se redefini\'o de forma clara: maximizar el
retorno esperado del cliente, tratando la utilidad de FinPUC como una
consecuencia \emph{ex post} del saldo activo administrado, sujeto a una
comisi\'on mensual de 0,5\% sobre dicho saldo.

El an\'alisis permite concluir que la metodolog\'ia \textbf{Black-Litterman con
\textit{view} de momentum general calibrada} es la alternativa m\'as conveniente
para los perfiles \textbf{Arriesgado} y \textbf{Muy arriesgado}, ya que mejora
el retorno del cliente frente al caso base \textbf{Markowitz}: obtiene una
mejora de Score P4 de 12.4\% bajo la agregaci\'on mediana robusta, 85.7\% bajo
la agregaci\'on promedio y 5.2\% en Sharpe. Desde la perspectiva del cliente,
genera mayor riqueza terminal esperada; desde la de FinPUC, la mejora tambi\'en
es favorable porque un mayor saldo activo administrado se traduce en mayor
utilidad por la comisi\'on mensual, sin introducir comisiones por transacci\'on ni
por rebalanceo.

\subsection{Calibraci\'on independiente de \textit{views}}

La calibraci\'on secuencial en cinco etapas se realiz\'o de forma independiente
para cada \textit{view}, sobre tres horizontes estrictamente separados, evitando
contaminar el per\'iodo de evaluaci\'on. Cada metodolog\'ia llega a la
simulaci\'on con su propia configuraci\'on congelada (por ejemplo, momentum
general con $c=0.80$ y $\tau=0.01$; desempleo macro con $c=0.50$ y $\tau=0.20$),
de modo que la comparaci\'on final contrasta metodolog\'ias en su mejor versi\'on
y no depende de una selecci\'on prematura de \textit{view}.

\subsection{Comparaci\'on de \textit{views}}

Entre las cuatro familias evaluadas bajo el modelo \'unico de abandono, momentum
general domina en t\'erminos econ\'omicos con una mejora mediana de +12.4\% en
Score P4 frente a Markowitz base (+85.7\% en promedio) y +5.2\% en Sharpe. Esta
ventaja conlleva mayor turnover (34.6\%) y concentraci\'on sectorial (HHI
0.261). Desempleo macro ofrece la mejor estabilidad operacional (turnover
7.3\%, HHI 0.157) pero retrocede en Score P4 mediano ($-$10.5\%). Las
\textit{views} basadas en capitalizaci\'on de mercado muestran deterioro tanto
en Sharpe como en Score P4 mediano ($-$4.7\%). La recomendaci\'on final prioriza
momentum general como \textit{view} base, con desempleo macro como alternativa
defensiva para perfiles conservadores.

\subsection{Abandono y utilidad por saldo}

El abandono semanal es marginal una vez que el cliente est\'a invertido: la
retenci\'on se ubica en torno a 850--870 clientes activos de cada 1\,000 al final
del horizonte. La utilidad de FinPUC sigue directamente al saldo administrado, de
modo que las metodolog\'ias que generan mayor riqueza terminal y sostienen la
retenci\'on producen tambi\'en mayor utilidad (USD~412 para momentum general
frente a USD~251 de Markowitz base).

\subsection{Validaci\'on y robustez}

La validaci\'on rolling de tres horizontes muestra que Black-Litterman calibrado
agrega m\'as valor en perfiles agresivos y escenarios con pandemia, donde la
estimaci\'on hist\'orica de Markowitz es menos confiable. En perfiles
conservadores y sin pandemia, la mejora es marginal, lo que es consistente con
que el prior de equilibrio ya est\'a cerca de la soluci\'on \'optima cuando los
datos son estables. La separaci\'on estricta de horizontes
$\mathcal{H}_1\,|\,\mathcal{H}_2\,|\,\mathcal{H}_3$ es una fortaleza
metodol\'ogica central que deber\'ia mantenerse en cualquier iteraci\'on futura.
Adicionalmente, los panoramas proyectados de retorno muestran que momentum
general conserva ventaja agregada en el caso neutro y en el desfavorable, por lo
que la recomendaci\'on no depende de un \'unico retorno esperado puntual.

\subsection{Limitaciones}

Las principales limitaciones son: (a) los KPIs de retorno/volatilidad usados en
Monte Carlo son fijos y no se re-optimizan dentro de la simulaci\'on; (b) la
distribuci\'on normal de los retornos semanales subestima eventos extremos, lo
que probablemente produce estimaciones optimistas de riqueza terminal; (c) el
Score P4 por mediana reduce el efecto de outliers, pero no penaliza
directamente turnover ni concentraci\'on sectorial; (d) la activaci\'on inicial
del cliente y la comisi\'on de 0,5\% mensual son supuestos no calibrados con
datos reales; y (e) la \textit{view} de desempleo usa una tasa fija en lugar de
series macroecon\'omicas observadas.

\subsection{Trabajo futuro}

Las extensiones prioritarias son: reemplazar la normalidad por t-Student o
Johnson SU en Monte Carlo, incorporar \textit{drawdown} respecto al m\'aximo de
riqueza como gatillo de abandono, usar series macro observadas (BLS/FRED) para
la \textit{view} de desempleo, calibrar la activaci\'on y abandono con datos
reales de clientes, y penalizar la rotaci\'on en la optimizaci\'on para contener
el turnover de momentum general.

\section{Carta Gantt: avance y fases restantes}

La Tabla~\ref{tab:plan_trabajo_futuro} resume la Carta Gantt de trabajo
posterior a esta entrega, diferenciando avance ya completado y tareas pendientes
para la fase final del proyecto.

\begin{table*}[h]
\centering
\caption{Carta Gantt resumida: avance actual y fases restantes.}
\label{tab:plan_trabajo_futuro}
\scriptsize
\begin{tabular}{@{}p{0.18\textwidth}p{0.14\textwidth}p{0.24\textwidth}p{0.34\textwidth}@{}}
\toprule
Fase & Periodo & Avance actual & Trabajo futuro \\
\midrule
Calibraci\'on independiente y validaci\'on &
12 mayo -- 19 junio 2026 &
Calibraci\'on secuencial por \textit{view}, validaci\'on rolling de tres horizontes y comparaci\'on de las cuatro \textit{views} bajo un \'unico modelo de abandono ya implementados. &
Profundizar la comparaci\'on con nuevas distribuciones de retorno y validar la sensibilidad del ranking final. \\
\addlinespace

Modelo de negocio por saldo administrado &
12 mayo -- 19 junio 2026 &
Comisi\'on de 0,5\% mensual sobre saldo activo y utilidad por retenci\'on implementadas y simuladas. &
Re-ejecutar la evaluaci\'on econ\'omica con el dataset completo y validar la convergencia de Monte Carlo. \\
\addlinespace

Mejoras del modelo &
19 -- 26 junio 2026 &
Se incorpor\'o Score P4 mediana como m\'etrica robusta y panoramas de retorno como sensibilidad post-optimal. &
Incorporar distribuciones no normales, \textit{drawdown} como gatillo de abandono y penalizaci\'on de turnover. \\
\addlinespace

Cierre del proyecto &
13 -- 26 junio 2026 &
Informe 3 en redacci\'on y principales resultados ya consolidados. &
Preparar la presentaci\'on final, documentar el modelo definitivo, ordenar el repositorio GitHub y realizar correcciones finales antes de la entrega. \\
\bottomrule
\end{tabular}
\end{table*}
